% 中文优化LaTeX模板 - 专门解决中文识别问题
% 智慧水利教材专用

\documentclass[UTF8,12pt,a4paper,oneside]{ctexbook}

% ==== 中文字体配置 ====
% 确保中文字体正确设置
\setCJKmainfont{SimSun}[
    BoldFont=SimHei,
    ItalicFont=KaiTi
]
\setCJKsansfont{SimHei}
\setCJKmonofont{FangSong}

% 英文字体配置
\setmainfont{Times New Roman}
\setsansfont{Arial}
\setmonofont{Consolas}

% ==== 基础包 ====
\usepackage{graphicx}
\usepackage{xcolor}
\usepackage{amsmath}
\usepackage{longtable}
\usepackage{booktabs}
\usepackage{array}

% ==== 页面设置 ====
\usepackage[
    top=2.5cm,
    bottom=2.5cm,
    left=3cm,
    right=2.5cm
]{geometry}

% 行距
\usepackage{setspace}
\onehalfspacing

% ==== 超链接 ====
\usepackage{hyperref}
\hypersetup{
    colorlinks=true,
    linkcolor=blue,
    urlcolor=blue,
    citecolor=blue,
    bookmarksopen=true,
    pdfstartview=FitH,
    unicode=true  % 重要：支持Unicode字符
}

% ==== 代码高亮 ====
\usepackage{listings}
\usepackage{xcolor}

% 定义颜色
\definecolor{codegreen}{rgb}{0,0.6,0}
\definecolor{codegray}{rgb}{0.5,0.5,0.5}
\definecolor{codepurple}{rgb}{0.58,0,0.82}
\definecolor{backcolour}{rgb}{0.95,0.95,0.92}

% 基本代码样式
\lstdefinestyle{mystyle}{
    backgroundcolor=\color{backcolour},   
    commentstyle=\color{codegreen},
    keywordstyle=\color{blue},
    numberstyle=\tiny\color{codegray},
    stringstyle=\color{codepurple},
    basicstyle=\ttfamily\footnotesize,
    breakatwhitespace=false,         
    breaklines=true,                 
    captionpos=b,                    
    keepspaces=true,                 
    numbers=left,                    
    numbersep=5pt,                  
    showspaces=false,                
    showstringspaces=false,
    showtabs=false,                  
    tabsize=2,
    frame=single,
    rulecolor=\color{black},
    extendedchars=false,  % 避免扩展字符问题
    inputencoding=utf8    % 明确指定UTF-8输入
}

\lstset{style=mystyle}

% 特定语言设置
\lstdefinestyle{javascript}{
    style=mystyle,
    language=JavaScript,
    morekeywords={const, let, var, function, async, await}
}

\lstdefinestyle{python}{
    style=mystyle,
    language=Python,
    morekeywords={def, class, import, from, as}
}

\lstdefinestyle{csharp}{
    style=mystyle,
    language=[Sharp]C,
    morekeywords={var, async, await, public, private}
}

% ==== 图片设置 ====
\graphicspath{{images/}}
\setkeys{Gin}{width=0.8\textwidth}

% ==== 标题设置 ====
\ctexset{
    chapter = {
        format = \Large\bfseries\centering\color{blue},
        beforeskip = 20pt,
        afterskip = 30pt
    },
    section = {
        format = \large\bfseries\color{blue},
        beforeskip = 15pt,
        afterskip = 15pt
    },
    subsection = {
        format = \normalsize\bfseries,
        beforeskip = 12pt,
        afterskip = 12pt
    }
}

% ==== 文档信息 ====
\title{\textbf{智慧水利平台架构与开发}}
\author{教材编写组}
\date{\today}

% ==== 文档开始 ====
\begin{document}

% 标题页
\maketitle
\newpage

% 目录
\tableofcontents
\newpage

% 正文内容
\section{智慧水利平台架构与开发}\label{ux667aux6167ux6c34ux5229ux5e73ux53f0ux67b6ux6784ux4e0eux5f00ux53d1}

\subsection{第一章 概述}\label{ux7b2cux4e00ux7ae0-ux6982ux8ff0}

\subsubsection{1.1 基本概念}\label{ux57faux672cux6982ux5ff5}

\textbf{智慧水利}是运用物联网、云计算、大数据、人工智能等现代信息技术，对水利工程进行智能化管理和运维的新模式。

主要特点包括：

\begin{enumerate}
\def\labelenumi{\arabic{enumi}.}
\tightlist
\item
  \textbf{数据驱动}：基于海量数据分析决策
\item
  \textbf{智能预警}：实时监测水情变化
\item
  \textbf{精准调度}：优化水资源配置
\item
  \textbf{协同管理}：多部门信息共享
\end{enumerate}

\subsubsection{1.2 技术架构}\label{ux6280ux672fux67b6ux6784}

系统采用\textbf{分层架构}设计：

\begin{itemize}
\tightlist
\item
  \textbf{感知层}：传感器网络、监测设备
\item
  \textbf{网络层}：通信传输、数据传递\\
\item
  \textbf{数据层}：数据存储、处理分析
\item
  \textbf{应用层}：业务功能、用户界面
\end{itemize}

\begin{Shaded}
\begin{Highlighting}[]
\CommentTok{\# 数据采集示例代码}
\KeywordTok{class}\NormalTok{ WaterDataCollector:}
    \KeywordTok{def} \FunctionTok{\_\_init\_\_}\NormalTok{(}\VariableTok{self}\NormalTok{, station\_id):}
        \VariableTok{self}\NormalTok{.station\_id }\OperatorTok{=}\NormalTok{ station\_id}
        \VariableTok{self}\NormalTok{.sensors }\OperatorTok{=}\NormalTok{ []}
    
    \KeywordTok{def}\NormalTok{ collect\_data(}\VariableTok{self}\NormalTok{):}
        \CommentTok{"""采集水位、流量数据"""}
\NormalTok{        data }\OperatorTok{=}\NormalTok{ \{}
            \StringTok{\textquotesingle{}时间\textquotesingle{}}\NormalTok{: datetime.now(),}
            \StringTok{\textquotesingle{}水位\textquotesingle{}}\NormalTok{: }\VariableTok{self}\NormalTok{.get\_water\_level(),}
            \StringTok{\textquotesingle{}流量\textquotesingle{}}\NormalTok{: }\VariableTok{self}\NormalTok{.get\_flow\_rate(),}
            \StringTok{\textquotesingle{}温度\textquotesingle{}}\NormalTok{: }\VariableTok{self}\NormalTok{.get\_temperature()}
\NormalTok{        \}}
        \ControlFlowTok{return}\NormalTok{ data}
    
    \KeywordTok{def}\NormalTok{ get\_water\_level(}\VariableTok{self}\NormalTok{):}
        \CommentTok{\# 获取水位数据}
        \ControlFlowTok{return} \FloatTok{125.6}  \CommentTok{\# 单位：米}
\end{Highlighting}
\end{Shaded}

\subsubsection{1.3 发展趋势}\label{ux53d1ux5c55ux8d8bux52bf}

智慧水利发展呈现以下趋势：

\begin{quote}
\textbf{重要提示}：数字化转型是水利现代化的必由之路。
\end{quote}

\textbf{发展方向}： - 🌊 全流域一体化管理 - 🤖 人工智能深度应用\\
- 📱 移动端智能服务 - 🔗 区块链技术应用

\subsection{第二章
系统设计}\label{ux7b2cux4e8cux7ae0-ux7cfbux7edfux8bbeux8ba1}

\subsubsection{2.1 需求分析}\label{ux9700ux6c42ux5206ux6790}

根据水利部门实际需求，系统需要实现：

\begin{enumerate}
\def\labelenumi{\arabic{enumi}.}
\tightlist
\item
  \textbf{实时监测}：24小时不间断监控
\item
  \textbf{预警预报}：提前发现风险隐患
\item
  \textbf{决策支持}：辅助管理决策
\item
  \textbf{应急响应}：快速处置突发事件
\end{enumerate}

\subsubsection{2.2 数据库设计}\label{ux6570ux636eux5e93ux8bbeux8ba1}

核心数据表包括：

\begin{longtable}[]{@{}lll@{}}
\toprule\noalign{}
表名 & 说明 & 主要字段 \\
\midrule\noalign{}
\endhead
\bottomrule\noalign{}
\endlastfoot
stations & 监测站点 & 站点ID、名称、位置、类型 \\
sensors & 传感器信息 & 设备ID、型号、参数、状态 \\
water\_data & 水文数据 & 时间、水位、流量、降雨量 \\
alerts & 预警信息 & 预警级别、内容、处理状态 \\
\end{longtable}

\begin{Shaded}
\begin{Highlighting}[]
\CommentTok{{-}{-} 创建水文数据表}
\KeywordTok{CREATE} \KeywordTok{TABLE}\NormalTok{ water\_data (}
    \KeywordTok{id}\NormalTok{ BIGINT }\KeywordTok{PRIMARY} \KeywordTok{KEY}\NormalTok{,}
\NormalTok{    station\_id }\DataTypeTok{VARCHAR}\NormalTok{(}\DecValTok{50}\NormalTok{) }\KeywordTok{NOT} \KeywordTok{NULL} \KeywordTok{COMMENT} \StringTok{\textquotesingle{}站点ID\textquotesingle{}}\NormalTok{,}
\NormalTok{    measure\_time DATETIME }\KeywordTok{NOT} \KeywordTok{NULL} \KeywordTok{COMMENT} \StringTok{\textquotesingle{}监测时间\textquotesingle{}}\NormalTok{,}
\NormalTok{    water\_level }\DataTypeTok{DECIMAL}\NormalTok{(}\DecValTok{10}\NormalTok{,}\DecValTok{2}\NormalTok{) }\KeywordTok{COMMENT} \StringTok{\textquotesingle{}水位(米)\textquotesingle{}}\NormalTok{,}
\NormalTok{    flow\_rate }\DataTypeTok{DECIMAL}\NormalTok{(}\DecValTok{15}\NormalTok{,}\DecValTok{2}\NormalTok{) }\KeywordTok{COMMENT} \StringTok{\textquotesingle{}流量(立方米/秒)\textquotesingle{}}\NormalTok{,}
\NormalTok{    rainfall }\DataTypeTok{DECIMAL}\NormalTok{(}\DecValTok{10}\NormalTok{,}\DecValTok{2}\NormalTok{) }\KeywordTok{COMMENT} \StringTok{\textquotesingle{}降雨量(毫米)\textquotesingle{}}\NormalTok{,}
\NormalTok{    temperature }\DataTypeTok{DECIMAL}\NormalTok{(}\DecValTok{5}\NormalTok{,}\DecValTok{2}\NormalTok{) }\KeywordTok{COMMENT} \StringTok{\textquotesingle{}水温(摄氏度)\textquotesingle{}}\NormalTok{,}
\NormalTok{    created\_time DATETIME }\KeywordTok{DEFAULT} \FunctionTok{CURRENT\_TIMESTAMP}
\NormalTok{);}
\end{Highlighting}
\end{Shaded}

\subsection{小结}\label{ux5c0fux7ed3}

本章介绍了智慧水利的基本概念、技术架构和发展趋势，为后续章节的深入学习奠定了基础。

\begin{center}\rule{0.5\linewidth}{0.5pt}\end{center}

\emph{注：本文档用于测试中文LaTeX转换效果}

\end{document}